\section{The denominator detected by residues on the whole fibre}
\label{sec:residue}

In the application to \eqref{eq:relative-local-cyclic}, the coefficient
below is $\beta=\alpha+t$.  Write $p_0\beta=a/m$ in lowest terms, with
$m>0$.  An $m$th root of a base uniformizer produces a log
pluricanonical generator with character $\zeta^{-a}$ for
$\zeta\in\mu_m$.  We compare its
$p/p_0$th power with the pullback of a degree-$p$ form on the whole
reduced fibre.  Divisible adjunction will first make their ratios
constant on each normalized component.  Residues along the conductor
will then make those constants equal, forcing the character to be
trivial and hence $m\mid p/p_0$.

\begin{lemma}[The reduced fibre and the cyclic residue test]
\label{lem:slc-residue-test}
Fix a finite set $\Phi\subset[0,1]\cap\Q$ and a positive integer
$p_0$ clearing its denominators.  There is a positive integer $p$,
depending only on $\Phi$ and $p_0$, divisible by $p_0$, with the following
property.  Let
\[
 f\colon N\longrightarrow C
\]
be a projective contraction from a normal integral variety to a smooth
complex curve, let $c\in C$, and suppose $1\leq\dim N\leq4$.  Set
$F=f^*[c]$ and $T=F_{\mathrm{red}}$.  Assume that $N$ is
$\Q$-factorial, that $(N,T+H)$ is dlt, and that $H$ is an effective
horizontal $\Q$-divisor with coefficients in $\Phi$.  Suppose, on a
neighbourhood of $F$, that there are $\beta\in\Q$ and
$\psi\in\C(N)^*$ such that the following is an equality of actual
$\Q$-divisors, for a fixed canonical divisor:
\begin{equation}
 \label{eq:slc-local-relation}
 p_0(K_N+T+H)+\Div\psi
       =p_0\beta F.
\end{equation}
Then $p\beta\in\Z$.
\end{lemma}

\begin{proof}
We may shrink $C$ about $c$ throughout.  In particular, we take a
uniformizer $z\in\C(C)$ with $\Div_C z=[c]$ on this
neighbourhood.  The fibre $T$ is connected and projective, since $f$ is
a projective contraction.  Its irreducible components have dimension
$\dim N-1$: the divisor $F$ is an effective Cartier divisor on the
integral variety $N$, and its support is the whole fibre.  When
$\dim N=1$, the contraction $f$ is an isomorphism of normal curves.
Comparing the coefficient at $c$ in
\eqref{eq:slc-local-relation} gives $p_0\beta\in\Z$, so this case is
immediate.  Henceforth $1\leq\dim T\leq3$.

\smallskip
\noindent\emph{Adjunction on the entire reduced fibre.}
Write $T=\bigcup_i T_i$.  By dlt adjunction and the structure theorem for
dlt strata, each $T_i$ is normal and
\begin{equation}
 \label{eq:slc-component-adjunction}
 (K_N+T+H)|_{T_i}\sim_{\Q}K_{T_i}+\Theta_i,
\end{equation}
where $(T_i,\Theta_i)$ is lc and $\Theta_i$ is the effective different;
see \cite[Sections~4.1--4.2 and Theorems~4.16 and~4.19]{Kollar2013}.
We use the residue form of this adjunction
\cite[Definition~11.14, author version]{KollarFamilies2023}.  If $n$ is sufficiently
divisible, the $n$th tensor power of the ordinary residue at the generic
point of $T_i$ extends to an isomorphism
\[
 \OO_N\bigl(n(K_N+T+H)\bigr)|_{T_i}
   \simeq \OO_{T_i}\bigl(n(K_{T_i}+\Theta_i)\bigr).
\]
The different is defined so that the generic residue extends uniquely
to this canonical isomorphism once the ambient Cartier index is
cleared.  Here $T_i$ is normal, so its normalization does not change
the target.  The isomorphism therefore identifies actual rational
pluricanonical forms.  The same form of adjunction will be used for
normal dlt components upstairs.
Intersections of distinct components of $T$ are unions of lc strata.
At the generic point of an intersection of codimension two in $N$, the
pair is simple normal crossing, precisely two components of $T$ pass
through the point, and $H$ is absent.  The corresponding conductor
prime on either branch occurs in $\Theta_i$ with coefficient one.

The remaining coefficients of the differents belong to the set
\begin{equation}
 \label{eq:slc-different-coefficients}
 \mathcal D(\Phi)=
 \left\{
   \frac{r-1+b}{r}\ \middle|\
   \begin{gathered}
    r\in\N_{>0},\quad n_\phi\in\Z_{\geq0},\\
    b=\sum_{\phi\in\Phi\cup\{1\}}n_\phi\phi\leq1
   \end{gathered}
 \right\}\cap[0,1].
\end{equation}
This is the usual coefficient rule for the different; see
\cite[Lemma~4.1]{HMX2014}.  The positive elements of
$\Phi\cup\{1\}$ have a positive minimum, so the possible sums $b\leq1$
form a finite set.  For each such sum, the displayed coefficients are
nondecreasing in $r$ and converge to one.  Consequently $\mathcal D(\Phi)$ is a
fixed rational DCC set.  We include the conductor coefficient one in
this set.

To apply the index theorem to $T$, we must glue this component
adjunction on the reduced scheme itself.  We will prove that $T$ is
demi-normal and that divisible residues identify its log
pluricanonical sheaf with an ambient restriction.  The degree needed
for this adjunction may depend on the pair.  Conductor-compatible
descent is also the issue addressed by the admissible-section methods
of Fujino \cite[Definition~4.1 and Lemma~4.2]{Fujino2000} and Gongyo
\cite[Section~5]{Gongyo2013}; here we give the local comparison needed
for the later uniform degree.

Apply the depth theorem
\cite[Theorem~11.18, author version]{KollarFamilies2023} to the dlt pair
$(N,T+H)$ with the effective integral divisor
$D=T\leq\lfloor T+H\rfloor$ and $L=0$.  It gives that
$\OO_N$, $\OO_N(-T)$, and $\OO_T$ are Cohen--Macaulay.
Here $\OO_N(-T)$ is the ideal of the reduced union of the prime
divisors $T_i$: on the normal variety $N$, it is the intersection of
their height-one prime ideals, equivalently the regular functions
vanishing to order at least one at every $T_i$.  Thus the subscheme in
the depth theorem is the reduced fibre used here.  In particular, $T$
satisfies $S_2$.

The normality of the $T_i$ and the generic normal crossing description
show that $T$ has only smooth points and ordinary double crossings in
codimension one.  It is also seminormal.  By the square--cube
criterion, an element of the total quotient ring whose square and cube
are regular lies in the local ring at each codimension-one point,
where seminormality is explicit.  The $S_2$ extension property then
puts it in $\OO_T$.  Thus $T$ is demi-normal.  Its normalization is the
disjoint union of the $T_i$, and its conductor divisor is the union of
the double-crossing divisors described above.

Remove the coefficient-one conductor divisors from the $\Theta_i$ and
push the remaining boundary cycles to $T$; denote the result by $B_T$.
Off the conductor the normalization is an isomorphism in codimension
one, so this does not double count any prime.  The boundary $B_T$ is
effective, has coefficients in $\mathcal D(\Phi)$, and has no component
contained in the conductor.

Choose an open immersion $j:T^\circ\hookrightarrow T$ whose complement
has codimension at least two, such that $T^\circ$ has only smooth points
and ordinary double crossings.  We may also arrange that at each crossing
the ambient pair is simple normal crossing with just the two vertical
branches, and that the support of $B_T|_{T^\circ}$ is contained in the
smooth locus.  The dualizing
sheaf of $T^\circ$ is invertible, and the relevant multiples of
$B_T|_{T^\circ}$ are Cartier.  For such a multiple $n$, interpret
$\OO_T(n(K_T+B_T))$ initially as
\[
 j_*\bigl(\omega_{T^\circ}^{\otimes n}
       \otimes\OO_{T^\circ}(nB_T|_{T^\circ})\bigr).
\]
We claim that for every sufficiently divisible positive integer $n$,
the residue maps give a natural isomorphism
\begin{equation}
 \label{eq:slc-ambient-restriction}
 \OO_N\bigl(n(K_N+T+H)\bigr)|_T
       \simeq \OO_T\bigl(n(K_T+B_T)\bigr).
\end{equation}
Choose $n$ even and divisible enough that the ambient sheaf is
invertible and the component residue isomorphisms hold.  Away from the
conductor, these isomorphisms give
\eqref{eq:slc-ambient-restriction} on $T^\circ$.  At a node, take local
normal crossing coordinates $x,y$ for the two branches.  Hypersurface
adjunction for $xy=0$ identifies the restriction of the ambient log
canonical line with the dualizing line of the node.  The two iterated
residues of $dx/x\wedge dy/y$ differ by a sign; their $n$th tensor
powers give the dualizing-sheaf gluing and have equal next residues
when $n$ is even.  This proves the claimed residue identification on
$T^\circ$.

An invertible sheaf on the $S_2$ scheme $T$ equals the extension of
its restriction from $T^\circ$.  Applying this to the left side of
\eqref{eq:slc-ambient-restriction} and using the definition of the
right side proves the isomorphism on $T$.  In particular, the defined
extension is invertible in these degrees, and the isomorphism retains
the actual rational residue forms on every normalized branch.

It follows that $K_T+B_T$ is $\Q$-Cartier.  Its pullback to the
normalization is $K_{T_i}+\Theta_i$, with the conductor included, so
$(T,B_T)$ is slc.  Equation~\eqref{eq:slc-local-relation}, together with
$F=\Div(f^*z)$, makes $K_N+T+H$ rationally linearly trivial
near $T$.  Restricting a sufficiently divisible trivialization in
\eqref{eq:slc-ambient-restriction} gives
\begin{equation}
 \label{eq:slc-fibre-trivial}
 K_T+B_T\sim_{\Q}0.
\end{equation}

\smallskip
\noindent\emph{A uniform degree on the slc fibre.}
Apply global ACC to the projective normalized component pairs
$(T_i,\Theta_i)$, whose adjoints are numerically trivial by
\eqref{eq:slc-fibre-trivial}.  Their dimensions are at most three and
their coefficients lie in the fixed DCC set
$\mathcal D(\Phi)$; hence their nonzero coefficients lie in a fixed
finite rational set \cite[Theorem~1.5]{HMX2014}.  The same is true of
$B_T$.  The bounded index theorem for projective slc log Calabi--Yau
pairs in dimensions at most three now gives a common integer
annihilating $K_T+B_T$ \cite[Corollary~1.6, arXiv version~1]{JiangLiu2021}.
Enlarging this integer, choose once and for all an even multiple $p$ of
$p_0$, depending only on $\Phi$ and $p_0$, such that
\begin{equation}
 \label{eq:slc-uniform-index}
 pB_T\text{ is integral},\qquad p(K_T+B_T)\sim0.
\end{equation}
Choose a nowhere-vanishing generator $u$ of this trivial log
pluricanonical line.  On the normalization it is a tuple of rational
$p$-pluricanonical forms $u_i$ on the $T_i$, with the prescribed boundary
poles and conductor gluing.

The comparison uses two degrees for different purposes.  A sufficiently
divisible degree $n=bp$ makes $u^b$ a generating ambient restriction
and will control divisors of componentwise ratios.  The conductor
comparison will then take place in the uniform degree $p$, using only
the ordinary node description there.

\smallskip
\noindent\emph{The root cover.}
Write
\begin{equation}
 \label{eq:slc-denominator}
 p_0\beta=\frac{a}{m},\qquad
 a\in\Z,\quad m\in\N_{>0},\quad\gcd(a,m)=1.
\end{equation}
Thus the goal is $m\mid p/p_0$. The case $m=1$ is immediate; assume $m>1$.  Let $C'\to C$ be the normalization in
$\C(C)(w)$, where $w^m=z$, and let $\pi\colon Y\to N$ be the
normalization of the dominating component of $N\times_C C'$.  The
polynomial $w^m-z$ is Eisenstein at $c$.  Thus $C'$ has a unique point
$c'$ over $c$, with total ramification of degree $m$ there.  Since $f$
is a contraction of normal varieties, $\C(C)$ is algebraically closed
in $\C(N)$; in characteristic zero this is a regular function-field
extension.  It is consequently linearly disjoint from
$\C(C')/\C(C)$.  The map $\pi$ has degree $m$ and cyclic Galois group
$\mu_m$, acting by $w\mapsto\zeta w$.
Write $f_Y:Y\to C'$ for the induced morphism and put
$T_Y=(\pi^{-1}T)_{\mathrm{red}}$. Figure~\ref{fig:root-cover}
records these maps and their reduced special fibres; the total space
$Y$ is the normalization of the base change.

\begin{figure}[ht]
\centering
\begin{tikzcd}[row sep=large,column sep=huge]
 T_Y \arrow[r] \arrow[d,hook] & T \arrow[d,hook] \\
 Y \arrow[r,"\pi"] \arrow[d,"f_Y"'] & N \arrow[d,"f"] \\
 C' \arrow[r,"w^m=z"'] & C
\end{tikzcd}
\caption{The normalized cyclic base change and its whole reduced special
fibres over $c'$ and $c$. The covering group may permute components of
$T_Y$, so the residue comparison must use their conductor gluing.}
\label{fig:root-cover}
\end{figure}

The morphism $f_Y$ is projective.  Its function
field is regular over $\C(C')$, because regularity is preserved by this
base change.  Stein factorization therefore gives
$(f_Y)_*\OO_Y=\OO_{C'}$. In particular, $T_Y$ is connected and
projective.
The divisor $p_0(K_N+T+H)$ in
\eqref{eq:slc-local-relation} is integral.  If
$F=\sum_i e_iT_i$, comparison of coefficients in that equation and
\eqref{eq:slc-denominator} gives $m\mid e_i$ for every $i$.
At the generic point of $T_i$, write $f^*z=\tau^{e_i}v$ with $v$ a unit.
After dividing $w$ by $\tau^{e_i/m}$, the normalized extension is
obtained by extracting an $m$th root of $v$ and is \'{e}tale.  Away from
$T$ the function $f^*z$ is a unit, and the same conclusion holds.
Hence $\pi$ is \'{e}tale in codimension one.

Put $H_Y=\pi^*H$, and choose canonical divisors compatibly.  Since
$\pi$ is quasi-\'{e}tale and the coefficient-one divisor pulls back
without divisorial ramification, we have
\begin{equation}
 \label{eq:slc-cover-crepant}
 K_Y+T_Y+H_Y=\pi^*(K_N+T+H).
\end{equation}
The pair on the left is dlt.  Indeed, finite crepant discrepancy
comparison multiplies a downstairs log discrepancy by the
ramification index of a divisorial prolongation
\cite[Proposition~5.20]{KollarMori1998}.  Thus the upstairs pair is lc,
and every upstairs lc centre maps onto a downstairs lc centre.
The pair downstairs is simple normal crossing near the generic point
of each such centre.  A finite quasi-\'{e}tale map is \'{e}tale over the
smooth locus of its target, by purity of the branch locus
\cite[Tag~0BMB]{StacksProject}.  Consequently the upstairs pair is
simple normal crossing at the generic points of all its lc centres.
A log resolution chosen to be an isomorphism over this log smooth open
has no exceptional divisor of log discrepancy zero, which is the dlt
condition.  This argument does not require $Y$ to be $\Q$-factorial.
In particular, each irreducible component $T_{Y,j}$ of $T_Y$ is normal.

Let $\theta$ be the rational top differential form defining $K_N$.
Define the rational $p_0$-pluricanonical form
\begin{equation}
 \label{eq:slc-cover-generator}
 s=(\pi^*\theta)^{\otimes p_0}\,\pi^*\psi\,w^{-a}.
\end{equation}
Since $\pi^*F=m\Div_Y w$, equations
\eqref{eq:slc-local-relation}--\eqref{eq:slc-cover-crepant} give
\[
 \Div(s)+p_0(T_Y+H_Y)=0.
\]
Thus $s$ generates the degree-$p_0$ divisorial log pluricanonical sheaf
on $Y$.  In particular its sufficiently divisible powers generate the
corresponding invertible log pluricanonical sheaves.

\smallskip
\noindent\emph{Residue comparison on each component.}
Let $T_i$ be the component of $T$ dominated by $T_{Y,j}$, and write
$\rho_j\colon T_{Y,j}\to T_i$ for the induced finite morphism.  At their
generic points the map is \'{e}tale.  The ordinary generic residue
\[
 v_j=\res_{T_{Y,j}}\bigl(s^{p/p_0}\bigr)
\]
is therefore a nonzero rational $p$-pluricanonical form.  The form
$\rho_j^*u_i$ is also nonzero, and the quotient
\begin{equation}
 \label{eq:slc-residue-ratio}
 r_j=\frac{v_j}{\rho_j^*u_i}\in\C(T_{Y,j})^*
\end{equation}
is a rational function.  We prove that $r_j$ is constant by comparing
sufficiently divisible powers.

Choose a multiple $n=bp$ sufficiently divisible for
\eqref{eq:slc-ambient-restriction} and for component adjunction upstairs.
Then $v_j^b$ is a nowhere-vanishing generator of the degree-$n$
adjunction sheaf on $T_{Y,j}$.  The same is true of $\rho_j^*u_i^b$.
For the latter assertion, work near any point of $T$.  By
\eqref{eq:slc-ambient-restriction}, the section $u^b$ is a generator of
the restriction of the invertible sheaf
$\OO_N(n(K_N+T+H))$.  Relative to a local ambient generator it is a unit
on $T$.  Such a unit lifts, after restricting the neighbourhood, to a
unit on $N$.  Hence $u^b$ is the residue restriction of a local ambient
log generator.  Pull that generator back to $Y$.  It remains a log
generator by \eqref{eq:slc-cover-crepant}; its residue on $T_{Y,j}$
generates the upstairs adjunction sheaf.  At the generic point residue
commutes with the \'{e}tale pullback, so this residue is exactly the
rational form $\rho_j^*u_i^b$.  This identifies the rational forms
everywhere they are defined, and proves the assertion at every point
of $T_{Y,j}$.

It follows that $r_j^b$ has zero divisor on $T_{Y,j}$, and hence so does
$r_j$.  A rational function with zero divisor on a normal projective
integral variety is an invertible global function, and therefore a
nonzero complex constant.  We have proved
\begin{equation}
 \label{eq:slc-component-scalar}
 v_j=r_j\rho_j^*u_i,\qquad r_j\in\C^*.
\end{equation}

\smallskip
\noindent\emph{Gluing the scalars and killing the character.}
Suppose two distinct components $T_{Y,j}$ and $T_{Y,j'}$ meet.  Their
intersection contains a stratum of codimension two in $Y$ along which
the pair is generically simple normal crossing.  Its image is a
codimension-two lc centre downstairs.  At the generic point of this
image the map $\pi$ is \'{e}tale and the downstairs pair is simple
normal crossing.  The two downstairs vertical branches are distinct:
\'{e}tale inverse images of one smooth divisor cannot give two
intersecting branches.  Moreover $H$ is absent at this generic
crossing.

At the generic point $\xi$ of the upstairs crossing, choose the two
normal coordinates $x,y$ and a regular ambient $(\dim Y-2)$-form
$\eta$ whose restriction is a nonzero top form on the smooth stratum.
Such choices are available in normal crossing coordinates.  Since
$s^{p/p_0}$ is a log generator, it has the local form
\[
 s^{p/p_0}
   =\varepsilon
      \left(\frac{dx}{x}\wedge\frac{dy}{y}\wedge\eta\right)^{\otimes p},
 \qquad \varepsilon\in\OO_{Y,\xi}^*.
\]
The two next residues differ by $(-1)^p$ and are nonzero because
$\varepsilon$ is a unit.  They therefore agree.  The components of
the tuple $u$ have the same matching, nonzero next residues downstairs:
$u$ is a generator of the slc log pluricanonical line, whose local
description at the node is the even tensor power of the dualizing
line.  Their \'{e}tale pullbacks retain this property.  Taking next
residues in \eqref{eq:slc-component-scalar} now gives $r_j=r_{j'}$.

Since $T_Y$ is connected, its component-incidence graph is connected.
All the constants $r_j$ are consequently equal to one scalar
$r\in\C^*$.  In tuple notation on the normalized components of $T_Y$,
we have
\begin{equation}
 \label{eq:slc-global-scalar}
 \bigl(\res_{T_{Y,j}}(s^{p/p_0})\bigr)_j
     =r\,\pi^*u.
\end{equation}
The tuple $\pi^*u$ is invariant under the covering group, with the
natural action also permuting the components.  On the other hand,
\eqref{eq:slc-cover-generator} shows that for $\zeta\in\mu_m$,
\[
 \zeta^*(s^{p/p_0})=\zeta^{-ap/p_0}s^{p/p_0}.
\]
Residue is natural under this action, including its permutation of the
components.  The nonzero scalar comparison
\eqref{eq:slc-global-scalar} forces this character to be trivial.
Thus $m\mid ap/p_0$.  Since $\gcd(a,m)=1$, we obtain
$m\mid p/p_0$.  Finally
\[
 p\beta=\frac{p}{p_0}\frac{a}{m}\in\Z,
\]
which proves the lemma.
\end{proof}

Apply Lemma~\ref{lem:slc-residue-test} to
\eqref{eq:relative-local-cyclic}, with $\beta=\alpha+t$ and the fixed
coefficient set $\Phi\cup\{1\}$.  It gives a single multiple $p$ of $p_0$
for which \eqref{eq:relative-coefficient-target} holds at every prime
$P$ of the determination $W$.  Consequently $pM_W$ is integral and
Cartier.  Since the moduli b-divisor descends on $W$, the b-divisor
$p\mathbf M$ is b-Cartier.  All constructions used only the dimensions
at most four and the fixed coefficient set.  This completes the proof
of Theorem~\ref{thm:relative-denominators}.
